\documentclass[11pt,letterpaper]{article}
\usepackage{technical-implementation}
% Shared non-semantic helper definitions; one semantic module is loaded by each root.
\newcommand{\PTSource}[1]{\par\smallskip{\footnotesize\color{Meta}\textbf{Book crosswalk:} #1}\par\smallskip}
\newcommand{\PTOutcome}[1]{\begin{tipbox}[Competency outcome]#1\end{tipbox}}
\newcommand{\PTEvidence}[1]{\begin{notebox}[Required evidence]#1\end{notebox}}
\newcount\ptlinecount
\newcommand{\PTLines}[1]{\par\ptlinecount=0\loop\ifnum\ptlinecount<#1\par\noindent\rule{\linewidth}{0.25pt}\vspace{0.22cm}\advance\ptlinecount by1\repeat\par}
\newcommand{\PTField}[1]{\noindent\textbf{#1}\quad\rule{0.66\linewidth}{0.3pt}\par\vspace{0.22cm}}
\newcommand{\PTRating}{\noindent\textbf{Rating:}\quad 0\quad 1\quad 2\quad 3\quad 4\qquad\textbf{Evidence ref:}\rule{0.30\linewidth}{0.3pt}\par\vspace{0.14cm}}
\newcommand{\PTTask}[2]{\par\medskip\noindent\textbf{Task #1.} #2\par}

\setDocTitle{Chapter 9: Deterministic Bottom-Up Parsing}
\setDocSubtitle{Foundations, reasoning, guided experiments and assessment}
\setDocVersion{1.0}\setDocStatus{Engineering education / source-aligned}
\setDocOwner{Compiler Engineering}
\setDocDate{30 October 2026}
\setDocSummary{Chapter 9 learning companion.}
\begin{document}\makedoctitle\makedocsummary
\tableofcontents\clearpage
\section{Chapter orientation}
This guide follows Chapter 9 of Grune and Jacobs, \emph{Parsing Techniques: A Practical Guide}, second edition. Printed pages 263--342. The focus is on reconstructing the chapter's reasoning, verifying explicit assumptions, and translating its central idea into a controlled engineering experiment. This is an independently authored learning guide, not a substitute for the book.
\PTSource{Chapter 9, Deterministic Bottom-Up Parsing; printed pages 263--342.}
\section{Learning objectives}
\begin{enumerate}
\item Define and differentiate operator precedence using a small grammar and its accepted language.
\item Define and differentiate lr(0) item sets using a small grammar and its accepted language.
\item Construct an example or counterexample illustrating canonical lr(1), slr(1), lalr(1).
\item Evaluate the operational constraints of shift/reduce and reduce/reduce diagnostics and retain reproducible evidence.
\end{enumerate}
\section{Core topics and review questions}
\subsection{Operator precedence}
Write the exact input domain, grammar assumptions, data structures and observable outputs. Distinguish a mathematical language property from the behavior of one implementation.
\textbf{Self-test:} Which invariant or decision criterion makes this topic precise, and what input would falsify an incorrect implementation?
\subsection{LR(0) item sets}
Construct two contrasting examples. Give a complete trace of states or sentential forms and indicate the condition under which the process halts.
\textbf{Self-test:} Which invariant or decision criterion makes this topic precise, and what input would falsify an incorrect implementation?
\subsection{Canonical LR(1), SLR(1), LALR(1)}
Identify a counterexample, ambiguity, cycle or boundary condition; explain whether it is a grammar property, an algorithm limitation, or an implementation bug.
\textbf{Self-test:} Which invariant or decision criterion makes this topic precise, and what input would falsify an incorrect implementation?
\subsection{Shift/reduce and reduce/reduce diagnostics}
State the cost model, failure modes and evidence needed for a defensible design choice. Do not extrapolate results beyond the grammar family studied.
\textbf{Self-test:} Which invariant or decision criterion makes this topic precise, and what input would falsify an incorrect implementation?
\section{Technical deep dive: LR items and control states}
An LR(0) item identifies a right-hand side and dot position. For a state $I$, $\operatorname{closure}(I)$ adds items $B\to\bullet\gamma$ whenever an item expects nonterminal $B$ after the dot, repeating to a fixed point. $\operatorname{goto}(I,X)$ advances the dot across symbol $X$ and closes the result. LR(1) items further carry lookahead; LALR state merging can introduce conflicts absent in canonical LR(1). A shift/reduce conflict must be tied to an input/context witness or a reasoned precedence policy.
\section{Worked reasoning prompt}
Construct closure and goto on the augmented grammar $S\prime\to S$, $S\to a$. Derive each intermediate state by hand; write the invariant that justifies advancing to the next state. Identify the accepted result and one structurally similar rejected result.
\section{Implementation and laboratory assignment}
Generate canonical item sets and record table conflicts for an ambiguous arithmetic grammar. Record the grammar, input corpus, expected outputs, maximum input size and a regression test for the principal failure mode. For each test, independently justify the expected result rather than reading it from the program.
\section{Failure analysis and technical discussion}
Do not claim SLR/LALR are equivalent to canonical LR(1) for all grammars. Explain how your implementation would detect the failure and which diagnostics can distinguish grammar invalidity from resource exhaustion.
\section{Evidence checklist}
\begin{itemize}
\item Grammar and accepted-language specification with version and fixtures.
\item Hand-executed trace and implementation trace for the same accepted and rejected cases.
\item At least one adversarial or boundary input and an explanation of the expected behavior.
\item Correctness argument, complexity claim with assumptions, and unit-test output.
\item Cross-reference to a related chapter and one explicit unresolved question.
\end{itemize}
\section{Independent examination}
\begin{enumerate}
\item Explain the algorithm without using implementation-language details.
\item Give a proof sketch or invariant for acceptance and termination on the specified grammar class.
\item Change one grammar rule and predict which tests must change before running them.
\item Compare behavior on malformed input, ambiguous input, and input beyond configured resource limits.
\end{enumerate}

\end{document}
